\chapter{Titre du chapitre 4}

\section{Introduction} 

Texte Texte Texte Texte Texte Texte Texte Texte Texte Texte Texte Texte Texte Texte Texte Texte Texte Texte Texte Texte Texte Texte Texte Texte Texte Texte Texte Texte Texte Texte Texte.

\section{Listes}

\begin{itemize}
	\item Premier item ;
	\item Deuxième item ;
	\item Dernier item.
\end{itemize}

\begin{itemize}[itemsep=0pt]
	\renewcommand{\labelitemi}{\scriptsize$\blacksquare$}
	\item Premier item ;
	\item Deuxième item ;
	\item Dernier item.
\end{itemize}

\begin{enumerate}
	\item Premier item ;
	\item Deuxième item ;
	\item Dernier item.
\end{enumerate}

\begin{enumerate}[label=\alph*)]
	\item Premier item ;
	\item Deuxième item ;
	\item Dernier item.
\end{enumerate}

\section{Graphiques}
La figure \ref{fig:exemple1} est un exemple avec dimensions originales (scale=1).

\begin{figure}[H]
	\centering
	\includegraphics[scale=1]{image}
	\caption{Titre de la figure.}
	\label{fig:exemple1}
\end{figure}


\bigskip
Figure avec dimensions modifiées d'une façon arbitraire
\begin{figure}[H]
	\centering
	\includegraphics[width=5cm, height=5cm]{image}
	\caption{Titre de la figure.}
	\label{fig:exemple2}
\end{figure}

\bigskip
Figure avec deux graphiques
\begin{figure}[H]
	\begin{minipage}{.5\linewidth}
		\centering
		\includegraphics{image}
	\end{minipage} \hfill
	\begin{minipage}{.5\linewidth}
		\centering
		\includegraphics{image}
	\end{minipage}
	\caption{Titre de la figure.}
	\label{fig:exemple3}
\end{figure}

\bigskip
Deux figures l'une à coté de l'autre
\begin{figure}[H]
	\begin{minipage}[h]{.5\linewidth}
		\centering
		\includegraphics{image}
		\caption{Titre de la figure.}
	\end{minipage} \hfill
	\begin{minipage}[h]{.5\linewidth}
		\centering
		\includegraphics{image}
		\caption{Titre de la figure.}
		\label{fig:exemple4}
	\end{minipage}
\end{figure}


\section{Mathématiques}
{\noindent\bfseries Exemple 1 :} Équations simples

\noindent Soit une fonction $f$ définie par
\[
f(x) = 2x-1
\]
Soit
\begin{equation} \label{eq:1}
	f(x) = 0
\end{equation}
La solution de l'équation (\ref{eq:1}) est $x=1/2$.

\bigskip
{\noindent\bfseries Exemple 2 :} Équations alignées
\begingroup
\setlength{\jot}{15pt}
\begin{align}
	& y=e^{x+1} \\
	& y=2\sin{x}+1 \\
	& y=\sqrt{x+1} \\
	& y=\left( \frac{1 + x}{2 + x^2} \right)^{2}\\
	& s=\sum_{i=1}^{n} x_{i}^{2} \\
	& p=\prod_{i=1}^{n} x_{i}^{2} \\
	& I=\int_{0}^{\pi} \sin x \, dx \\
	& l= \lim _{x\rightarrow \infty} \log(x+1)/x
\end{align}
\endgroup


\bigskip
{\noindent\bfseries Exemple 3 :} Équation en plusieurs lignes
\begin{align*}
	f(x)&=(x+1)(x+2)^2 \\
	&=(x+1)(x^2+2x+4) \\
	&=x^3+3x^2+6x+4
\end{align*}

\bigskip
{\noindent\bfseries Exemple 4 :} Équations avec parenthèse
\begin{equation}
	{\begin{array}{rl}
			\text{Minimiser} &
			\left\{
			\begin{array}{l}
				f_1(x)=1-\exp\bigl(-\sum_{i=1}^n(x_i-1/\sqrt{n})\bigr), \\
				f_2(x)=1-\exp\bigl(-\sum_{i=1}^n(x_i+1/\sqrt{n})\bigr),
			\end{array}
			\right.
	\end{array}}
\end{equation}

\bigskip
{\noindent\bfseries Exemple 5 :} Équation avec plusieurs cas
\begin{equation}
	\begin{array}{l}
		v_{j,t}^i=
		\left\{
		\begin{array}{ll}
			V_{max,j} & \text{si }v_{j,t}^i>V_{max,j} \\[0.2cm]
			-V_{max,j} & \text{si }v_{j,t}^i<-V_{max,j} \\[0.2cm]
			v_{j,t}^i & \text{sinon}
		\end{array}
		\right.
	\end{array}
\end{equation}

\bigskip
{\noindent\bfseries Exemple 6 :} Matrices

Soient $\mathbf{A}=[a_{ij}]$ une matrice $n\times p$ et $\mathbf{B}=[b_{ij}]$ une matrice $p\times q$. Le produit de ces deux matrices est la matrice $\mathbf{C}=\mathbf{A}\mathbf{B}$, qui a pour dimension $n\times q$ et s'écrit:

\begin{displaymath}
	\mathbf{C}=[c_{ij}] \ \text{avec} \ c_{ij}=\sum_{k=1}^{p}a_{ik}b_{kj}, \ \text{pour}\ 1\leq i\leq n \ \text{et} \ 1\leq j\leq q
\end{displaymath}

\bigskip
\noindent Soient \\
\begin{displaymath}
	\mathbf{A}=
	\begin{pmatrix}
		-2 & 1 \\
		0 & 1 \\
		-1 & 2
	\end{pmatrix} \quad \text{et} \quad
	\mathbf{B}=
	\begin{pmatrix}
		6 & 0 & 2\\
		3 & -2 & 2
	\end{pmatrix}
\end{displaymath}\\
On a donc
\begin{displaymath}
	\mathbf{C}=
	\begin{pmatrix}
		-2 & 1 \\
		0 & 1 \\
		-1 & 2
	\end{pmatrix} \cdot
	\begin{pmatrix}
		6 & 0 & 2\\
		3 & -2 & 2
	\end{pmatrix}=
	\begin{pmatrix}
		-9 & -2 & -2 \\
		3 & -2 & 2 \\
		0 & -4 & 2
	\end{pmatrix}
\end{displaymath}

\section{Tableaux}
Utilisez l'outil en ligne : latex-tables.com 

\begin{table}[h]
	\caption{Exemple d'un tableau simple.}
	\centering
	\begin{tabular}{|l|c|r|}
		\hline
		Colonne 1 & Colonne 2 & Colonne 3 \\ \hline
		1.1 & 1.2 & 1.3 \\ \hline
		2.1 & 2.2 & 2.3 \\
		\hline
	\end{tabular}
\end{table}

\begin{table}[h]
	\setlength{\arrayrulewidth}{0.5mm}
	\setlength{\tabcolsep}{18pt}
	\renewcommand{\arraystretch}{1.5}
	\caption{Exemple d'un tableau avec des espacements ajustés.}
	\centering
	\begin{tabular}{|l|c|r|}
		\hline
		Colonne 1 & Colonne 2 & Colonne 3 \\ \hline
		1.1 & 1.2 & 1.3 \\ \hline
		2.1 & 2.2 & 2.3 \\ \hline
	\end{tabular}
\end{table}

\begin{table}[h]
	\caption{Exemple d'un tableau avec des largeurs personnalisées.}
	\centering
	\begin{tabular}{|>{\raggedright\arraybackslash}p{4cm}|>{\centering\arraybackslash}p{3cm}|>{\raggedleft\arraybackslash}p{4cm}|}
		\hline
		Colonne 1 & Colonne 2 & Colonne 3 \\ \hline
		1.1 & 1.2 & 1.3 \\ \hline
		2.1 & 2.2 & 2.3 \\ \hline
	\end{tabular}
\end{table}

\begin{table}[h]
	\caption{Exemple d'un tableau sans les traits verticaux.}
	\centering
	\begin{tabular}{lcr}
		\toprule
		Colonne 1 & Colonne 2 & Colonne 3 \\ \midrule
		1.1 & 1.2 & 1.3 \\
		2.1 & 2.2 & 2.3 \\ \bottomrule
	\end{tabular}
\end{table}

\begin{table}[h]
	\caption{Exemple d'un tableau avec une colonne commune.}
	\centering
	\begin{tabular}{|l|c|r|}
		\hline
		Colonne 1 & \multicolumn{2}{c|}{Colonnes 2 et 3} \\ \hline
		1.1 & 1.2 & 1.3 \\ \hline
		2.1 & 2.2 & 2.3 \\ \hline
	\end{tabular}
\end{table}

\begin{table}[h]
	\caption{Exemple d'un tableau avec une ligne commune.}
	\centering
	\begin{tabular}{|l|c|r|}
		\hline
		Colonne 1 & Colonne 2 & Colonne 3 \\ \hline
		\multirow{2}{*}{12.1} & 1.2 & 1.3 \\ \cline{2-3}
		& 2.2 & 2.3 \\ \hline
	\end{tabular}
\end{table}

\begin{table}[h]
	\caption{Exemple d'un tableau avec plusieurs cellules communes.}
	\centering
	\begin{tabular}{|l|c|r|}
		\hline
		Colonne 1 & \multicolumn{2}{c|}{Colonnes 2 et 3} \\ \hline
		\multirow{2}{*}{12.1} & 1.2 & 1.3 \\ \cline{2-3}
		& 2.2 & 2.3 \\ \hline
		3.1 & 3.2 & 3.3 \\ \hline
	\end{tabular}
\end{table}

\begin{table}[h]
	\caption{Exemple d'un tableau avec un trait diagonal.}
	\centering
	\begin{tabular}{|l|c|c|}\hline
		\diagbox[width=10em]{Lignes}{Colonnes} &
		Colonne 1 & Colonne 2 \\ \hline
		Ligne 1 & 1.1 & 1.2 \\ \hline
		Ligne 2 & 2.1 & 2.2 \\ \hline
	\end{tabular}
\end{table}

\begin{table}[h]
	\caption{Exemple d'un tableau contenant une image et une liste.}
	\centering
	\begin{tabular}{ | c | m{5cm} | m{5cm} | }
		\hline
		Colonne 1 & Colonne 2 & Colonne 3 \\ \hline
		\begin{minipage}{.3\textwidth}
			\centering
			\vspace{0.5cm}
			\includegraphics[scale=0.7]{image}
			\vspace{0.5cm}
		\end{minipage}
		&
		%\begin{minipage}[t]{5cm}
		\begin{itemize}[leftmargin=*,itemsep=0pt]
			\item Premier item
			\item Deuxième item
			\item Dernier item
		\end{itemize}
		%\end{minipage}
		& 1.3
		\\ \hline
	\end{tabular}
\end{table}
